\chapter{Critical Literature Review}
\label{ch:literature}

\section*{Multimodal Monitoring in Robotic Welding}
\addcontentsline{toc}{section}{Multimodal Monitoring in Robotic Welding}
During robotic welding, a camera can observe changes in the weld pool while a microphone records the sound of the process. The two observations can respond differently to the same change. The main research question asks how to design and evaluate a process-aware architecture that uses such evidence to detect anomalies reliably. To judge the value of a second input, the joint model and its strongest single-input alternative must predict the same outcome on the same welds. Two primary welding studies illustrate why this comparison matters.

The dynamic multimodal measurement fusion network for welding (DMAW) combined image and audio Transformer representations to classify four welding states. Accuracy reached 96.1\% with fusion, compared with 90.2\% for the image model and 84.4\% for the audio model. The reported evaluation leaves the grouping of windows by welding run and performance during a controlled outage unresolved \citep[PDF p.~7, Table~III]{P003}. In a weld-penetration study, the image-only, audio-only, and fused models each reached 100\% accuracy on the clean test set. The fused model achieved 93.33\% under the reported interference condition. Determining the benefit of fusion during interference requires the single-input models to be evaluated under the corresponding conditions \citep[PDF pp.~18, 22, Tables~1--3]{Hua2026Weld}. Together these studies establish that learned visual--acoustic fusion already exists in welding and that its observed advantage depends on the endpoint and comparator.

The original proposal concerned Intel Welding. The subsequent audit documented previous use of the proposed test population, which prevented an independent weld-quality assessment (Chapter~\ref{ch:introduction}). The review therefore follows the design question into adjacent tasks. Industrial condition studies clarify how a physical measurement becomes a prediction target. Robotic and machining studies show what can be compared when a camera and machine signals observe the same operation. The four subquestions concern the representation of each input (SQ1), the value of combining inputs (SQ2), behavior when an input fails or degrades (SQ3), and the evaluation of uncertainty, decisions, and practical use (SQ4). Section~\ref{sec:literature-study-gaps} identifies the specific comparison each of the five implementations could make.

\section{Review Scope and Search Methodology}

The review traced industrial fusion claims to the methods, comparison tables, and results in the publications inspected. A gap here means a comparison that those results leave unresolved for a specified target, set of inputs, and assessment population. The search and access records limit how broadly that inference can be made.

\subsection{Research Themes and Search Questions}

The representation theme asks which observations are available, what transformations are applied, and whether a representation preserves information relevant to the target. The fusion theme asks whether additional inputs improve a common prediction task against credible single-input references. The robustness theme asks which absence, corruption, or distribution-change regime was actually evaluated. The evaluation theme examines independence, thresholds, uncertainty, resource measurement, and reproducibility. These questions deliberately include negative outcomes: a fusion mechanism should be informative even when an added modality provides little benefit.

The review also distinguishes two objectives often grouped together. Learning with several modalities to deploy a single-input model is a privileged-training problem. Combining several live inputs during deployment is a multimodal-inference problem. Cross-modal distillation and reconstruction connect these objectives, but their acquisition assumptions differ. The taxonomy of \citet{Baltrusaitis2019Multimodal} provides an organizing vocabulary; industrial claims are appraised in the original studies rather than inferred from that survey.

\subsection{Sources, Search Terms, and Search Period}
The literature review developed throughout the project. One documented search round focused on work published in 2016--2026, with selected earlier methodological contributions. Of the 120 candidates screened in that round, 39 met the inclusion criteria. The corresponding full-text appraisal covered 37 of those works. The later expansion identified \ExpandedPDFCount{} distinct full-text files after reconciling duplicate downloads and checking work versions. The current comparison register comprises \ReviewWorkCount{} distinct work records, rather than an estimate of the size or prevalence of the entire field.

Searches paired industrial monitoring and manufacturing terms with fusion stage, attention, missing inputs, noise, and domain shift. Dataset searches included MVTec AD, MVTec 3D-AD, Eyecandies, bearing datasets, and machine-sound benchmarks. Publisher and repository searches located publications. Crossref supplied metadata for identity checks \citep{Crossref2020REST}, and OpenAlex supported citation tracing \citep{Priem2022OpenAlex}. Zotero linked bibliographic records to the inspected PDFs and appraisal notes \citep{Zotero2026}. Appendix~\ref{app:literature} distinguishes the original screening population from the later full-text appraisal.

\subsection{Inclusion and Exclusion Criteria}
Direct industrial work concerns faults, condition, manufacturing quality, inspection or robotic task outcomes. Dataset papers and credible unimodal methods are also relevant because fusion needs a defined target and a strong reference. Transferable studies are labelled explicitly: optimization failures, missing-input mechanisms, uncertainty weighting and evaluation theory can inform industrial design without constituting industrial validation \citep{Wang2020What,P016,P023,Bengio2004CVVariance}.

A PDF can remain in the register while serving only as background or carrying major methodological qualifications. Simulated robot regression, RGB--geometry inspection, physical camera--telemetry monitoring and one-camera feature fusion occupy different strata \citep{Altinses2023Fuzzy,P026,Thoduka2024Handover,P035}. Inaccessible evidence is not presumed irrelevant, and missing reporting is not proof that an operation was omitted. Author terminology alone does not determine whether a paper supplies multimodal acquisition, inference-time robustness or a fair unimodal comparison.

\subsection{Evidence Extraction and Critical Appraisal}
Each appraisal records source identity and read version, target and prediction unit, physical inputs, transformations, fusion computation, controls, partitioning, training access and reliability interventions. Results are linked to pages, tables or figures. Decisive PDF tables were checked where extraction could obscure input columns, values or footnotes. Full-text appraisal establishes inspection, not experimental reproduction or endorsement of every author conclusion.

The comparison asks what was held constant and what changed. A common assessment cohort supports a pipeline comparison; it does not isolate fusion if encoders, supervision, augmentation or capacity also differ. Input-subset models, complete inference absence, corrupted present data and target-domain adaptation are coded separately. The comparison matrix in Appendix~\ref{app:literature}, Table~\ref{tab:literature-matrix}, supplies the work-level reference. Its compact codes require the accompanying protocol descriptions and source locators.

\subsection{Limitations of the Review}
The sources are neither exhaustive nor statistically independent. Many inspection papers reuse MVTec 3D-AD or Eyecandies, so repeated benchmark gains do not constitute repeated validation across factories. Preprints are version-specific and some publication details remain provisional. When few studies evaluated a condition, the review offers limited evidence about that condition.

The primary MVTec 3D-AD table gives 249 normal and 948 anomalous tests, totaling 1,197; some later accounts give different counts \citep[Table~1]{Bergmann2022The}. The welding penetration paper has inconsistent acquisition/training counts, and a lightweight inspection paper gives different NwfDD values in its principal and ablation tables \citep[PDF pp.~12, 15]{Hua2026Weld}\citep[PDF p.~10, Tables~IV, VI, VII]{P035}. These problems reduce the precision of particular comparisons without invalidating every observation in a source.

\section{Industrial Monitoring Tasks and Supervision}

Different monitoring tasks impose different assumptions about training data, observations, and acceptable error. Their scores are therefore not interchangeable evidence for a general architecture. The target definition must precede the comparison of representation or fusion mechanisms.

\subsection{Anomaly Detection, Defect Classification, and Failure Classification}
Normal-only anomaly detection learns a nominal reference, whereas supervised defect classification uses labelled fault or quality classes. MVTec AD and MVTec 3D-AD exemplify the first setting; welding penetration recognition exemplifies the second \citep{Bergmann2019MVTec,Bergmann2022The,Hua2026Weld}. HANDOVER introduces a different endpoint: whether a completed interaction succeeded or failed through no approach, no grasp, no release or object drop \citep{Thoduka2024Handover}. A task failure need not denote physical damage or identify its cause.

Supervision is consequently part of the comparison. The weakly supervised inspection study reintegrates MVTec data because training needs anomalous image-level labels, changing the ordinary normal-only contract \citep[PDF p.~8]{P035}. ZSG-IAD trains reports and masks before evaluating held-out categories or datasets; its zero-shot terminology concerns that boundary, rather than absence of anomaly-supervised training \citep[PDF pp.~6, 8--9]{P025}. Scores from these tasks cannot be combined as interchangeable evidence for a general architecture.

\subsection{Condition Estimation and Regression}

Continuous prediction requires attention to units, reference measurements, and the horizon at which a target is defined. Remaining useful life is not equivalent to an instantaneous defect probability. \citet{AlDulaimi2019A} evaluate a hybrid CNN--LSTM on simulated C-MAPSS trajectories using RMSE and an asymmetric prognostic score. Their two branches process the same selected telemetry channels; this is evidence about complementary representations of sensor data, not physical camera--sensor complementarity.

The synthetic robot benchmarks of \citet{Altinses2025Performance} address cross-modal prediction, reconstruction, and downstream regression. Errors for different reconstructed outputs cannot automatically establish which physical modality is most useful for one common target. This distinction matters for proxy supervision: a system may predict a derived condition index accurately without demonstrating that the index measures deterioration. A defensible regression comparison specifies the target construction, scale, fitted preprocessing, and common evaluation cases before interpreting numerical improvement as better condition monitoring.

\subsection{Dataset-Provided Labels, Derived Targets, and Proxy Supervision}

Reference labels have different evidential roles. Post-build CT provides flaw localization for the voxel-level monitoring study of \citet{Petrich2021Multimodal}; microscopy and density measurements support quality classes in \citet{Li2022A}; operando X-ray imaging guides process-state labels in \citet{Pandiyan2022Deep}. The reference instrument may therefore supervise the model without being available at inference. Counting it as an input modality would overstate the deployed system's information.

Derived labels require a further separation between prediction and independent validation. If a target is calculated from the same measurements supplied to the model, low prediction error may reflect recovery of that calculation rather than an independent physical state. Likewise, region-level annotations may assign a common quality class to many neighboring observations. These choices can be useful, but their spatial or temporal dependence affects the evaluation unit. The thesis consequently needs an explicit account of which targets are native annotations, independent measurements, or analytical proxies, and which conclusions each can support.

\subsection{Retrospective Assessment and Online Monitoring}
A predictor evaluated after a complete window is a retrospective assessor unless prediction-time availability is constrained. HANDOVER explicitly states that its baselines classify outcomes after the complete trial; causal online prediction remains future work \citep[PDF p.~7, Section~VI]{Thoduka2024Handover}. A post-pass photograph, layer summary and complete action therefore provide different information horizons from prospective monitoring \citep{P045,Petrich2021Multimodal}. This distinction does not invalidate offline classification; it limits the operational claim.

Computation must also be separated from observation. DMAW uses twenty frames at 37 frames per second while reporting a shorter forward latency. Approximately half a second of observation is therefore distinct from computational delay \citep[PDF pp.~6--7]{P003}. Likewise, varying CNN--LSTM input duration does not by itself establish warning before a defect occurs \citep{Pandiyan2022Deep}. An online claim needs a causal information cutoff, acquisition delay and defined end-to-end workload.

\section{Multimodal Data in Industrial Applications}

Complementarity depends on independent observations, not simply on distinct tensor shapes. The review therefore distinguishes physical acquisition channels from derived views, and distinguishes ordinary pairing from the stronger correspondence required by cross-modal reconstruction or localized process monitoring.

\subsection{Visual and Video Observations}
RGB, thermal images, process videos and geometry have different acquisition assumptions. Depth, point clouds reconstructed from depth and estimated normals may be several representations of one geometric source. GPAD combines RGB--depth features with point-cloud priors; this does not establish three independent physical sensors \citep[PDF pp.~4--7]{P024}. RGB--geometry benchmarks describe appearance and shape, rather than controller dynamics or motor current \citep{Bergmann2022The,Bonfiglioli2023The}.

Physical camera-plus-measurement precedents include thermal imaging with vibration, welding images with sound, tool photographs with acoustic emission and robot video with force--torque and gripper state \citep{Huo2022Multisensor,Hua2026Weld,P045,Thoduka2024Handover}. Their units range from static inspection to process windows and complete trials. The tool-wear camera is used after a cutting pass, whereas preceding acoustic emission describes active cutting \citep[PDF pp.~19--22]{P045}. A sequence of images can also be processed without learning temporal order: pooled frame features and a video Transformer are different representations even when both consume multiple frames.

\subsection{Sensor Time Series and Process Measurements}

Machine measurements vary in physical origin and semantic meaning. Accelerometers observe vibration, pressure sensors observe fluid behaviour, and controller logs may describe commands, estimates, or tracking errors rather than independently measured physical quantities. The harmonic-drive study combines high-frequency vibration with servo information, whereas PGMTKD transfers pressure--vibration teacher knowledge to a pressure-only student \citep{Yang2024Fault,Li2025PGMTKD}. Their sensor availability and deployment objectives differ despite both addressing missing information.

Representations can further multiply apparent modalities without adding measurements. \citet{Kullu2022A} combine current and vibration with raw and STFT views. \citet{Wan2024Pipeline} derive statistical ``vibration'' features from two acoustic sensors; they do not acquire an accelerometer channel. \citet{Wang2025A} combine source-domain features and decisions from vibration data. These studies remain useful architectural comparators, but neither a second transform nor a second domain establishes an additional physical observation. Acquisition inventories must therefore precede counting modalities or attributing gains to sensor diversity.

\subsection{Audio Observations}

Audio can capture machine or process events that are weakly visible, but the mapping between recorded sound and fault state depends on background conditions, microphone placement, and synchronization. MIMII deliberately mixes factory background recordings at multiple SNR levels; its fan baseline AUC decreases from .94 at +6~dB to .70 at -6~dB \citep[Table~3]{Purohit2019MIMII}. ToyADMOS provides miniature-machine recordings and an explicit noise-mixing protocol, making it a useful acoustic reference rather than a demonstrated camera--sensor advantage \citep{Koizumi2019ToyADMOS}.

In process monitoring, airborne acoustic signals may be paired with a melt-pool camera or a photodiode. The distinction is in the instruments. Signal-to-image encoding in \citet{Li2022A} preserves an acoustic--optical-signal experiment, whereas \citet{Xie2025Audiovisual} use actual camera observations. DCASE2020 Task2 provides audio-only baselines under an evaluation protocol that differs from those of later domain-shift challenges \citep{Koizumi2020Description}. Acoustic coverage in the review therefore motivates candidate representations without implying predictive audio evaluation in every thesis case study.

\subsection{Synchronization, Correspondence, and Modality Availability}

Fusion assumes that observations refer to a compatible event, region, or time interval. Stronger correspondence is required when a model predicts one modality from another or multiplies aligned discrepancy maps. RGB--point-cloud feature mapping exploits registered benchmark geometry; process fusion may instead need time synchronization and projection into a common spatial grid \citep{Costanzino2024Multimodal,Petrich2021Multimodal}. The digital-twin study registers sensor features to robot TCP positions, illustrating that position can support correspondence without serving as an ordinary predictive feature \citep{Chen2023Multisensor}.

Availability has several levels: a file may exist, a sample may have a paired observation, and a usable signal may still be absent or degraded. None implies the others. Modality completeness should therefore be recorded per prediction unit with an explicit pairing rule. The acoustic--photodiode study explicitly lacks modality registration, which qualifies conclusions drawn from its encoded-window fusion \citep{Li2022A}. The same distinction matters for camera and machine-sensor data. A pairing record can identify unmatched events even when both files are present.

\section{Representation Learning for Industrial Monitoring --- SQ1}

SQ1 concerns the relationship between representations and task requirements. The evidence supports neither a universal preference for learned features nor an assumption that a large pretrained encoder is always necessary. Representation quality interacts with supervision, preprocessing, sample size, and the resources allocated to fitting.

\subsection{Engineered Features and Conventional Machine Learning}
Engineered statistics, spectral descriptors and conventional classifiers remain important references. Their relevance follows from predictive performance and information access, not from architectural novelty. Classical geometric features can be highly competitive for 3D inspection, and a multi-channel neural model can lose to a selected sensor \citep{Horwitz2023Back,Goodarzi2023Deep}. A strong reference prevents a weak learned branch from defining an artificially easy fusion comparison.

Industrial vision also has lightweight alternatives. The weakly supervised inspection network uses one camera source, internal anomaly-aware features and a compact CNN; its reported cost is 2.33 million parameters and 5.65~ms per image \citep[PDF p.~10, Table~V]{P035}. This is relevant representation and cost evidence, but its anomalous training labels and modified MVTec supervision differ from normal-only inspection. Conventional and learned references should therefore be compared under the same target, assessment cases and fitted preprocessing rather than ranked by model family.

\subsection{Pretrained Visual Representations}
Pretrained CNN and Transformer features reduce industrial fitting requirements while importing substantial external learning. PatchCore, asymmetric student--teacher models and cross-modal feature mapping show useful nominal-data representations under different protocols \citep{Roth2022Towards,Rudolph2023Asymmetric,Costanzino2024Multimodal}. Their success supports transfer as a candidate strategy, not an architecture-independent guarantee.

Vision-language inspection imports semantic context as well. DSE-AD combines historical/current rail images with normal/abnormal position prompts and evaluates pixel localization on another rail domain \citep{P031}. Its paired-reference and prompt assumptions differ from isolated image classification. Similarly, a detector checkpoint can supply only convolutional features: that use does not imply execution of a Transformer encoder, decoder or detection head. A representation description must identify the actual forward path, trainable components and observation horizon before an improvement can be attributed to pretraining or attention.

\subsection{Temporal and Time-Series Representations}
Temporal structure can be represented through statistics, spectra, recurrence, attention or aggregation. PatchTST tokenizes channels into patches and shares Transformer weights across channel-independent paths. Its supervised, masked-pretraining and transfer experiments are different fitting regimes \citep[PDF pp.~3--8]{Nie2023PatchTST}. On Weather horizon 96, the patch-plus-channel-independent model reports MSE .152 against .164 without patching; this supports the design within forecasting, not superiority for robotic failure classification \citep[PDF p.~9, Table~7]{Nie2023PatchTST}.

Video-time attention addresses a different structure. TimeSformer separates temporal and spatial attention, while Video Swin introduces local spatiotemporal attention \citep{P034,P013}. Neither is equivalent to averaging independently encoded frames. The industrial CNN--LSTM RUL work also shows dependence on observation length within its simulated task \citep{AlDulaimi2019A}. SQ1 therefore requires explicit windows, ordering, normalization, pretraining and resource budgets. Selecting statistics after a bounded Transformer comparison remains a substantive answer to suitability.

\subsection{Audio Representations}
Acoustic encodings include MFCCs, log-mel spectra, wavelet maps, GASF transforms and learned features. The penetration network applies VMD and GASF, while DMAW adapts Transformer representations for welding observations \citep{Hua2026Weld,P003}. An image-shaped acoustic tensor remains sound evidence, not an independently acquired camera stream. Spectra must also use the same available interval as competing representations.

The Audio Spectrogram Transformer supplies a relevant pretrained candidate, but generic audio benchmarks do not establish industrial reliability or its advantage on every fault target \citep{P011}. MIMII and ToyADMOS make background composition and noise mixing explicit \citep{Purohit2019MIMII,Koizumi2019ToyADMOS}. This supports careful recording and normalization choices rather than an automatic architecture selection. The REASSEMBLE alignment checks prevented evaluation of an acoustic predictor (Chapter~6).

\subsection{Frozen Representations, Fine-Tuning, and Computational Constraints}

Freezing an encoder, fitting a small head, and adapting an entire network allocate very different statistical and computational resources. Comparing them requires disclosure of which parameters change and what data influence selection. Benchmark methods based on pretrained features can be effective with nominal training data, but their representation-learning cost is partly external to the industrial experiment \citep{Roth2022Towards,Costanzino2024Multimodal}. This is legitimate transfer learning, provided that the comparison does not present it as equivalent to training from scratch.

Distillation introduces another boundary between training and deployment. PGMTKD's multimodal teachers enrich a pressure-only student, while the gated acoustic--vibration study reports a distilled LightGBM student but leaves ambiguity about teacher-feature inputs and timing boundaries \citep{Li2025PGMTKD,Huang2026Rolling}. Resource claims must account for feature extraction, teachers required at inference, and preprocessing as well as the final predictor. For SQ1, the appropriate conclusion is a need to compare explicit fitting recipes under feasible budgets, not to assume that frozen features or full fine-tuning are intrinsically fairer.

\section{Multimodal Fusion Strategies --- SQ2}

Fusion is classified here by the location of information exchange. Early integration combines raw channels or engineered features; intermediate integration combines learned representations; output-level fusion combines predictions or anomaly scores. Hybrid systems use more than one stage. Author terminology is retained where needed, but it does not override this operational distinction.

\subsection{Unimodal References and Incremental Modality Value}
The strongest credible single-input model defines the central SQ2 reference. Welding makes the issue immediate: the penetration paper labels its own audio-only and image-only models as achieving 100\%, equal to clean fusion. Gains over other backbones cannot establish an added-modality effect \citep[PDF pp.~18, 22]{Hua2026Weld}. HANDOVER likewise reports 69.1\% for video plus force--torque and 67.9\% after adding gripper state within the I3D-D family \citep[PDF p.~7, Table~III]{Thoduka2024Handover}. These native-paper values are five-class outcome accuracies averaged over five trainings, not binary failure AP.

The tool-wear comparison reports MAE 6.6183 for its proposed model, 8.3002 for visual ResNet50 and 60.5118 for its acoustic reference; other multimodal architectures have errors above 45 \citep[PDF p.~24, Table~1]{P045}. Architecture differences prevent causal attribution to fusion alone, but the strong visual reference changes the interpretation. Transferable naive multimodal failures and industrial all-channel losses reinforce the need to report incremental effects against the best admissible component \citep{Wang2020What,Goodarzi2023Deep}.

\subsection{Feature Concatenation}
Concatenation combines representations without a complex allocation mechanism, making it a necessary reference. Its performance depends on feature scale, dimension, regularization and encoder quality. M3DM's direct-concatenation ablation is below its point-cloud reference, whereas the complete hybrid system improves on both inputs \citep[Table~4]{Wang2023Multimodal}. MADFlow provides another protocol-specific alignment-versus-concatenation comparison \citep{Li2025MADFlow}.

CFA supplies explicitly transferable evidence: for DLinear on Energy, first concatenation yields MSE .670 versus .263 without text and .243 with constrained fusion \citep[PDF p.~4, Table~4]{P023}. Yet its own constrained model also loses in some settings. Manufacturing fusion can similarly help one architecture and harm another under the same broad task \citep[Table~3]{Li2022A}. Concatenation should receive reasonable normalization and fitting; a weak implementation cannot establish the necessity of a more elaborate operator.

\subsection{Prediction-Level and Late Fusion}
Late fusion preserves separately fitted predictors and combines outputs through averaging, stacking or evidence rules. Thermal--vibration diagnosis uses separate CNN softmax outputs with weighted Dempster--Shafer aggregation. Its clean reported diagnostic rate is 100\%, compared with 99.06\% for IRT and 96.32\% for vibration \citep[PDF p.~7, Table~VII]{Huo2022Multisensor}. Fixed weights and the sample-CV protocol qualify the gain, but this is direct physical camera--sensor late-fusion evidence.

HANDOVER compares late outputs with intermediate sensor integration into RGB and optical-flow branches \citep[PDF pp.~5--7]{Thoduka2024Handover}. Inspection map fusion is different: cross-modal mapping and CPIR aggregate discrepancy maps that need not be probabilities \citep{Costanzino2024Multimodal,Shangguan2025CPIR}. Moreover, some papers call pre-classifier representation combination late fusion \citep{Wang2020What,Xie2025Audiovisual}. A learned prediction combiner therefore requires an explicit fitting bank, calibration contract and score scale, in addition to a common outer assessment.

\subsection{Gated and Context-Conditioned Fusion}
Gates can allocate embeddings, blend maps, alter prediction coefficients or balance training objectives. DMAW reports 96.1\% against 90.5\% for summation and 92.6\% for its gate comparator \citep[PDF p.~9, Table~V]{P003}. TC-MAF instead bounds auxiliary RGB evidence around a multimodal anchor using normal-training dispersion and local disagreement \citep[PDF pp.~3--5]{P026}. These are different adaptive mechanisms.

The conditioner is decisive. GCCF-AD describes dynamic fusion, but Eq.~8 specifies a learned scalar $\alpha=\sigma(\omega)$ without an explicit sample-quality input \citep[PDF p.~3]{P038}. PGMTKD's teacher allocation uses training losses, while task-loss balancing does not become a deployment quality measurement \citep{Li2025PGMTKD,P045}. GAF-Net uses progressive feature feedback, and trustworthy fusion uses evidence-based aggregation \citep{Huang2026Rolling,Ying2025Trustworthy}. These mechanisms differ in which weights can vary, which inputs determine them and whether the resulting predictions improve on a fixed combination. Process context does not itself impose physical laws.

\subsection{Comparability and Fairness of Fusion Evaluations}
A shared dataset name does not establish a matched experiment. Cohort, supervision, observation window, pretraining, feature dimensions, augmentation and selection exposure can explain gains. The tool-wear study changes architecture alongside fusion; incomplete federated imagery changes client data access as well as the integration mechanism \citep{P045,P051}. These are informative pipeline comparisons, but do not automatically isolate attention or quality conditioning.

Same-evidence controls support narrower attribution. TC-MAF retains its evidence sources while comparing bounded fusion with static mean, weighted scores, rank fusion and a learned gate \citep[PDF p.~7, Table~3]{P026}. Conversely, the disagreement-aware paper improves a local released-code baseline that is below the published M3DM result \citep[PDF p.~10, Table~4]{P033}. Distilled predictors also receive training information unavailable to ordinary unimodal models \citep{Li2025PGMTKD}. Such settings are legitimate when identified. Pipeline benefit, modality necessity and a causal fusion-operator effect remain different questions.

\section{Reliability and Robustness --- SQ3}

Robustness is defined by an intervention and a prediction setting. Missing channels, corrupted measurements, asynchronous observations, and a new machine are different problems. The matrix records evidence for each regime separately; general robustness language is not used to fill an untested cell.

\subsection{Missing Modalities}
Planned single-input deployment differs from unannounced loss in a model expecting complete inputs. PGMTKD's pressure-only student and incomplete RGB--geometry methods establish useful reduced-input inference, while ModDrop and missing-modality Transformers provide transferable mechanisms and counterexamples \citep{Li2025PGMTKD,Sui2026Incomplete,P005,P016}. They do not establish every possible outage pattern.

TC-MAF explicitly evaluates unavailable 3D: I-AUROC falls from .9788 to .9330 and P-AUPRO from .9901 to .9586 \citep[PDF p.~7, Table~3]{P026}. Federated substation work varies the fraction of incomplete clients, reaching 80\%, which differs from an arbitrary test-camera failure matrix \citep[PDF p.~29, Figure~8]{P051}. ZSG-IAD's subset-input rows do not fully specify shared-checkpoint and training exposure \citep{P025}. Missing training images in hot rolling constitute another regime \citep{Ma2025A}. Each comparison must state known masks, training availability, replacement policy and all-inputs-absent behavior.

\subsection{Noise and Partial Degradation}
Physical camera robustness studies already exist. Thermal--vibration diagnosis injects Gaussian noise and vibration bias/stuck faults; welding penetration recognition reports artificial interference; acoustic--thermal and acoustic--vibration studies examine further noise settings \citep{Huo2022Multisensor,Hua2026Weld,Ying2025Trustworthy,Huang2026Rolling}. These interventions differ in affected streams, severity and exposure.

At -2~dB, noisy IRT gives 95.26\%, rising to 97.6\% with fault-free vibration; noisy vibration gives 36.86\%, rising to 94.2\% with fault-free IRT \citep[PDF p.~7, Table~VIII]{Huo2022Multisensor}. The healthy companion is part of that condition. The welding paper's interference result is less reproducible because operator, magnitude and exposure matching are unclear \citep{Hua2026Weld}. Synthetic SNR tests do not establish natural simultaneous degradation, and ordinary defective objects are targets rather than corrupted modalities. Filtering or a harsh-environment title alone cannot fill a robustness cell.

\subsection{Temporal Misalignment and Distribution Shift}
A clock offset breaks correspondence; domain shift changes the assessment distribution. Multiple sampling rates or image repetition do not demonstrate offset tolerance. Registration studies establish why alignment matters, while OmniFuser's repeated within-pass images provide a strategy rather than a synchronization-failure benchmark \citep{Petrich2021Multimodal,Chen2023Multisensor,P021}.

Stronger transfer examples include training on one HANDOVER robot and assessing another, and Subway-trained rail localization evaluated on high-speed rail \citep[PDF p.~7, Table~V]{Thoduka2024Handover}\citep[PDF p.~10, Table~III]{P031}. Target-adaptive studies, target-validated condition studies and independent-build experiments have different access assumptions \citep{Wang2025A,Chu2024A,Huang2026Rolling,Bevans2023Heterogeneous}. Rebuilding normal memories on Eyecandies with fixed fusion hyperparameters is useful recipe transfer, but not untouched-target or cross-machine deployment \citep{P026,P033}. Genuine temporal offsets, calibration drift and new physical groups need separately identified protocols.

\subsection{Quality-Aware and Uncertainty-Aware Fusion}
Signal quality, confidence, calibrated probability and evidence consistency are different quantities. QMF motivates weighting through generalization upper bounds under assumptions about uncertainty, loss and a common distribution; this is not a guarantee of actual-risk superiority after arbitrary shift \citep[PDF pp.~4--5]{Zhang2023QMF}. The transferable UDML experiments train on controlled Gaussian intensities and acknowledge that sample difficulty can be mistaken for modality corruption \citep[PDF pp.~4--5, 9]{P022}.

QMF itself reports NYU scene-classification accuracy of 70.09\% on clean inputs, below concatenation at 70.30\%, but 55.60\% at corruption intensity ten, above late fusion at 51.99\% \citep[PDF p.~8, Table~1]{Zhang2023QMF}. These non-industrial RGB--depth results demonstrate a regime-specific benefit rather than universal dominance.

Results can also conflict across endpoints. The disagreement add-on improves local MVTec image AUROC from .779 to .800 but reduces Eyecandies image AUROC from .810 to .794 while improving pixel metrics \citep[PDF pp.~10--11, Tables~4--5]{P033}. TC-MAF's training-dispersion proxy has a smaller contribution than its evidence/fusion structure \citep{P026}. These findings require validation of the conditioning signal and outcomes; changing attention, an uncertainty plot or an evidence mass does not by itself establish calibrated reliability.

\subsection{Modality Dropout, Augmentation, and Fallback Mechanisms}
Training dropout exposes a model to selected absence patterns, imputation generates substitutes and deterministic fallback routes to an available predictor. These mechanisms can yield the same output under some conditions without learning the same recovery. Survivor routing should therefore be evaluated as a policy, while corrupted-but-present inputs test a different prediction problem.

Fuzzy regularization represents failures through partial corruption and whole-input replacement by constants or noise, evaluating Gazebo pose regression under random sample splits of overlapping windows \citep[PDF pp.~3--5]{Altinses2023Fuzzy}. This does not establish robustness to an unavailable input branch in a physical deployment. OmniFuser asserts dropout and masking without a quantified outage table \citep[PDF p.~10]{P021}. Generated replacements can harm: replacing clean depth with CycleGAN output worsens a synthetic robot benchmark's MSE from $2.1745\times10^{-3}$ to $13.472\times10^{-3}$ \citep[Table~7]{Altinses2025Performance}. Exposure-matched static controls, explicit replacements and surviving-input references are consequently necessary to distinguish learned recovery from imposed logic.

\section{Evaluation Beyond Accuracy --- SQ4}

Practical usefulness depends on how predictions are assessed, selected, and interpreted. Accuracy can be high when the target prevalence, grouping, or threshold makes an experiment easier than its intended deployment. SQ4 therefore requires a broader account of evidence than a single aggregate score.

\subsection{Grouped Evaluation and Leakage Prevention}
The independent unit may be a weld, tool, participant, recording, machine or build. Random windows can share acquisition context or raw samples. Fuzzy regularization forms twenty-step windows with stride one and randomly partitions examples; unseen-trajectory assessment is not established \citep[PDF pp.~4--5]{Altinses2023Fuzzy}. Weld-run grouping is also unresolved in the closest welding studies \citep{P003,Hua2026Weld}. These observations identify dependence risks without measuring an invented leakage magnitude.

HANDOVER supplies an explicit participant-disjoint official split, while an adapted implementation still needs participant information for safe inner folds \citep{Thoduka2024Handover}. Normalization, feature selection and HPO need the same separation \citep{P004,P017}. Out-of-fold stacking can protect outer assessment, but that does not prove every deeper meta-validation dependency is absent. Repeatedly choosing follow-ups after observing assessment outcomes also creates selection exposure even when each fit excludes its outer test labels.

\subsection{Discrimination, Calibration, and Decision Thresholds}
Ranking, probability quality and thresholded decisions answer different questions. AUROC does not specify a false-alarm operating point or calibrated probability. MVTec AD shows that good ranking can coexist with poor thresholded localization, while strict-FPR metrics reward another aspect of detection \citep{Bergmann2019MVTec,Costanzino2024Multimodal,Cheng2025Multimodal}.

The rail vision-language paper reports pixel metrics and optimal-threshold F1; these are not image-level alarm accuracy under a training-selected threshold \citep[PDF p.~7]{P031}. In manufacturing, accuracy .9921 alongside precision .2487 on an imbalanced layer illustrates the effect of prevalence and threshold \citep{Petrich2021Multimodal}. Validation-fitted calibration can improve in-distribution probability quality, while dataset shift can undermine confidence estimates \citep{Guo2017Calibration,P008}. Proper losses such as NLL and Brier assess probability forecasts, with an explicit distinction between loss orientation, calibration and discrimination \citep{P055}. A generated confidence field or evidence mass cannot replace probability-quality evaluation. Threshold selection belongs to admissible training/validation, while assessment reveals the resulting trade-off.

\subsection{Uncertainty and Statistical Comparisons}
Seed variation, acquisition-group sampling and selection uncertainty are distinct. HANDOVER's standard deviations describe five trainings on a fixed test set, whereas a group bootstrap of saved predictions samples assessment units conditional on fitted models \citep{Thoduka2024Handover}. Neither automatically includes all retraining or domain uncertainty.

CV folds share training data. The classical variance result establishes that no universally unbiased CV-variance estimator exists under its stated assumptions; ignoring covariance can substantially underestimate uncertainty \citep[PDF pp.~8--10]{Bengio2004CVVariance}. This does not prohibit useful conditional intervals. It requires an explicit estimand and dependence model. Small fusion effects, unclear repetition counts and contradictory significance prose need similar caution \citep{Gu2024Improved,Xu2025ModalityResilient}. Paired effects on common cases are preferable to unrelated score intervals. A zero-containing interval does not establish equivalence, and a positive point estimate does not establish practical importance without a meaningful criterion.

\subsection{Interpretability and Its Limits}
Scores, coefficients, gate traces and maps can expose model behavior without supplying a causal explanation. The attention debate includes both counterexamples to faithfulness and evidence that trained weights contain useful information; the appropriate conclusion is task-specific evaluation, rather than a title-level rule \citep{P036}. Allocation changes alone do not show why a model predicts correctly.

ZSG-IAD measures schema validity, language consistency and coarse-to-final-mask overlap. These support particular forms of self-consistency, not independent physical defect causation or calibrated confidence \citep{P025}. HANDOVER similarly leaves explanation through human-action predictions as future research \citep{Thoduka2024Handover}. In the thesis, observable quality responses and inspectable coefficients can support descriptive transparency. Expert understanding, causal attribution and usability require further evidence. Process context must also be distinguished from a physical equation constraint such as PGMTKD's pressure model \citep{Li2025PGMTKD}.

\subsection{Modularity, Reproducibility, and Computational Cost}
A modular pipeline enables controlled replacement, but does not itself establish robustness. An identified version, accessible code, executable processing and independently reproduced findings provide different assurances \citep{P009}. The local-versus-published M3DM discrepancy illustrates why dependencies, valid-point filtering, checkpoints and scoring conventions matter \citep{P033}.

Cost comparisons also need a common boundary. CPIR reports lower FPS than CFM, exposing a performance--speed trade-off \citep{Shangguan2025CPIR}. A lightweight image network's 5.65~ms excludes an independent sensor branch and acquisition delay \citep[PDF p.~10, Table~V]{P035}. TimeSformer's cost depends on ImageNet pretraining and workload \citep{P034}. Hardware, batch size, caching, synchronization and included stages determine what can be compared. Prospective deployment and operator usability remain separate validation questions rather than consequences of a successful offline forward pass.

\section{Critical Synthesis and Implications for This Thesis}
The strongest precedents for the welding question combine observations of the same physical operation and evaluate them against a defined outcome. RGB--geometry inspection, robot simulation, and studies of learning with missing inputs contribute evidence about particular components. Their results have different implications for the five implemented studies because the target, observation window, and conditions of deployment change.

\subsection{Findings Supported Across the Reviewed Studies}
Distinct physical observations can improve a task-specific pipeline. DMAW, thermal--vibration diagnosis, HANDOVER and tool-wear prediction provide relevant examples under their respective protocols \citep{P003,Huo2022Multisensor,Thoduka2024Handover,P045}. Controlled evidence combination can also improve RGB--geometry inspection, including a bounded rule with an explicit missing-3D result \citep{P026}. A broad claim that these combinations or reliability approaches are absent is therefore untenable.

The evidence does not select a universal winner. Planned single-input deployment, complete absence, synthetic degradation and target-domain adaptation remain different achievements \citep{Li2025PGMTKD,Altinses2023Fuzzy,P051}. Strong representations, labels and observation windows can matter as much as integration. The supported finding is conditional benefit under a defined task and protocol, with distinct limitations in acquisition independence, uncertainty and operational validation.

\subsection{Conflicting Findings and Unresolved Questions}
Negative findings are substantive evidence. Welding unimodal equality, lower HANDOVER performance after adding gripper state, worse multimodal tool-wear alternatives than a strong visual reference and Eyecandies image-score deterioration challenge unconditional fusion benefit \citep{Hua2026Weld,Thoduka2024Handover,P045,P033}. Category-specific geometry losses, harmful hydraulic-channel combinations and naive multimodal optimization failures reinforce that conclusion \citep{Horwitz2023Back,Bonfiglioli2023The,Goodarzi2023Deep,Wang2020What}.

Possible explanations include redundancy, optimization, noise, supervision and limited data; these studies do not identify one universal cause. Questions about independent groups, matched static controls, informative quality conditioning and calibrated risk remain protocol-specific. Where tables within a study disagree, the magnitude of the reported effect remains uncertain.

\subsection{Implications for Representation and Fusion Selection}
The welding results suggest an order for comparing architectures. First establish how well each physical input predicts the chosen target under a feasible fitting procedure (SQ1). Then compare equal averaging, a fitted static combination, and a context-dependent alternative with the strongest constituent predictor on the same assessment cases (SQ2). If the learned combination introduces another calibration step, a gain over a single branch estimates the benefit of the fitted pipeline. Isolating the contribution of the second input also requires a single-input control with the same opportunity for recalibration. If a gate is intended to handle failures, its static comparator needs comparable exposure to those failures during training (SQ3). The uncertainty of each contrast depends on the physical units that were independently assessed (SQ4).

\subsection{Implications for Experimental Design}
\label{sec:literature-study-gaps}
The following comparisons link the five implementations to the main question. Each starts from the closest published evidence and identifies what remains uncertain for the target actually available in that study.

\paragraph{PHM gear monitoring (SQ1, SQ2, and SQ4).}
Post-build computed tomography can provide an independent reference for locating flaws in manufacturing \citep{Petrich2021Multimodal}. In tool-wear prediction, an image taken after a pass and acoustic emission recorded during cutting have also been used together to estimate wear \citep[PDF pp.~19--24]{P045}. PHM poses a different measurement problem. Its gear photographs were collected periodically, while the challenge specifies sensor-only operation at deployment \citep{PHM2026Challenge}. The unresolved question is whether available sensor groups can predict a continuous condition proxy derived from post-run photographs across held-out gear experiments. Chapter~5 compares sensor representations and fusion of their embeddings with a training-mean reference (SQ1 and SQ2). Error in percentage points describes prediction of the photographic proxy (SQ4). Establishing its relation to physical tooth damage would require an independent measurement. Inference-time loss of a camera falls outside that comparison.

\paragraph{REASSEMBLE robotic assembly (SQ1--SQ4).}
The REASSEMBLE dataset paper benchmarks action segmentation and visual condition checks \citep[PDF pp.~10--13]{Sliwowski2025REASSEMBLE}. Published handover experiments already demonstrate multimodal classification of completed robotic trials \citep[Table~III]{Thoduka2024Handover}. For failed assembly actions, two comparisons remain open. Can a fitted video--sensor model rank failures better than a strong statistical sensor model? Does process-dependent weighting improve on equal averaging? Chapter~6 holds out complete recordings and compares visual, statistical-sensor, and time-series representations (SQ1), followed by feature concatenation, averaging, stacking, and gating on the same eligible actions (SQ2). Known absence and controlled corruption address different reasons to distrust an input (SQ3). Paired recording-level uncertainty, probability errors, and decisions at selected thresholds qualify the performance differences (SQ4). The observation window extends through action completion.

\paragraph{MSM milling (SQ1--SQ4).}
OmniFuser forecasts milling-tool wear states and cutting force \citep[PDF p.~11, Tables~IV--V]{P021}. A visual--acoustic study estimates tool wear after a cutting pass \citep[PDF p.~24, Table~1]{P045}. The MSM dataset instead identifies abnormal cutting intervals that warrant intervention. Its descriptor distinguishes these intervals from stable cutting with a defective tool and documents camera coverage in the laboratory \citep{Kim2026MSM}. The open comparison is whether video improves on machine telemetry for that interval-level distinction within cuts observed by both streams. Chapter~\ref{ch:msm} tests representations and static fusion against their own single-input predictors (SQ1 and SQ2), then compares performance under missing or degraded inputs (SQ3). Grouped evaluation and the potential disagreement between ranking and probability scores address SQ4. The covered laboratory dates bound the claim about new operating conditions.

\paragraph{ViFailback manipulation (SQ1, SQ2, and SQ4).}
\citet[PDF pp.~7--8, Tables~2--4]{Zeng2026ViFailback} use visual symbols to diagnose failures, evaluate responses to questions, and guide corrective actions. For the binary outcome of a complete manipulation episode on a held-out task, two comparisons remain open. Does adding a generic visual representation improve on robot-state measurements, and does transferring a visual encoder from REASSEMBLE change that comparison? Chapter~\ref{ch:vifailback} compares generic frozen features with those of the transferred encoder (SQ1) and assesses probability fusion against the constituent sensor model (SQ2). Separately reporting success and failure ranking and their macro-average makes the minority outcome visible (SQ4). The defined corruption procedures lacked a completed stress evaluation for the compared models, leaving the reliability question (SQ3) open for this task.

\paragraph{HANDOVER human--robot interaction (SQ1--SQ4).}
The published HANDOVER benchmark compares unimodal and fused models for five outcome classes. Within its I3D comparison, video plus force--torque signals achieved higher mean accuracy than the combination that also included gripper state \citep[PDF p.~7, Table~III]{Thoduka2024Handover}. The supplementary thesis study asks a narrower question. Over a common observation window, does averaging video and wrench predictions improve the ranking of binary handover failures over the better single-input model? Appendix~\ref{app:handover} compares the two representations and fixed averaging (SQ1 and SQ2), examines sensor noise, blur, and fallback to one input (SQ3), and assesses decisions for each outcome class (SQ4). The published participant-disjoint split supports the official test comparison. The missing participant-to-trial mapping prevented a separate participant-disjoint fit of the planned learned stacker.

\medskip
\noindent The five studies thus test different parts of the method proposed for welding. Their results can establish when a representation, fusion rule, or fallback policy helps on its own target. A welding-specific answer still requires an independently assessed weld-quality outcome, a defined decision time, and paired observations on welding sessions reserved for that assessment. Chapter~\ref{ch:discussion} interprets these comparisons.
